\chapter{Arabic logical traditions: syllogism, modality and signification}
\label{sec:p13-arabic-logical-traditions}

Arabic logical writing is not one unchanged system transmitted under a
different vocabulary. Al-Farabi, Avicenna and Averroes disagree about the
division of syllogisms and the interpretation and conversion of modal
propositions. Later commentators make Avicenna's formulations an object
of criticism, not merely a text to preserve. This account follows
Tony Street's \emph{Logic} in the \emph{Cambridge Companion to Arabic
Philosophy} and his study of the reception of \emph{Pointers} 1.6
\cite{Street2005LogicSupplement,Street2020PointersSignification}.
These are distinct presentations from his longer \emph{Arabic Logic}
chapter in the 2004 \emph{Handbook of the History of Logic}
\cite{Street2004ArabicLogicHHL}; no equivalence of their extents is assumed.

\section{A translated corpus and more than one logical project}
\label{sec:p13-arabic-corpus-traditions}

Street's shorter account concerns the Peripatetic logical tradition,
not all forms of reasoning practised in Arabic or in Muslim societies.
Its comparison of early and later writing is organized around the
problems inherited by Najm al-Din al-Katibi's \emph{Shamsiyya}. Ending
that investigation with the late thirteenth-century formation of a
textbook does not date the end of original logical activity.
\cite[pp.~247--249]{Street2005LogicSupplement}

The Syriac teaching tradition encountered by early Muslim scholars
included the \emph{Categories}, \emph{On Interpretation}, and the
assertoric opening of the \emph{Prior Analytics}. Translation into
Arabic did not proceed through one universally repeated sequence.
Summaries, individual treatises and commentaries were translated and
revised; Greek-to-Arabic work commonly used Syriac intermediaries.
The circles associated with al-Kindi and Hunayn ibn Ishaq differed in
their interests and resources. By about 900 the complete \emph{Organon}
was available in Arabic. Consequently the Latin distinction between
an old and a newly recovered logical corpus cannot simply become
the chronology of Arabic logic. Street also cautions against accepting
al-Farabi's account of a Christian prohibition on advanced logic as
an uncontested explanation of the earlier curriculum.
\cite[pp.~249--251]{Street2005LogicSupplement}

The Baghdad Peripatetics, including Abu Bishr Matta and al-Farabi,
worked toward close interpretation of Aristotle. The annotated Arabic
\emph{Organon} associated with their tradition made sustained exegesis
possible. Street's historical referral identifies Paris, Bibliotheque
nationale, arabe 2346 as a relevant witness; this is a referral to the
scholarship he cites, not a fresh collation here. The Farabian project
continued in Baghdad and in Andalusian work, including Averroes.
Avicenna, by contrast, took an independently defensible system to
govern how Aristotle should be read or corrected. Respect for a
predecessor need not mean adopting his inferential rules.
\cite[pp.~251, 263--264, notes~13--14 and 27--32]{Street2005LogicSupplement}
\index{al-Farabi}\index{Avicenna}\index{Averroes}

\section{Syllogistic organization and the matter of an argument}
\label{sec:p13-arabic-syllogistic-organization}

The \emph{Shamsiyya} organizes an introductory treatment around terms,
propositions and syllogisms. Terms include utterances, individual
concepts, universals and particulars, and definition. Propositions
include categorical and hypothetical forms, quantification, existence,
modality, contradiction, conversion and obversion. The syllogistic
includes mixed premises, hypothetical propositions, linked reasoning,
reduction, induction and example. Similar headings in Latin and
Arabic manuals do not establish identical systems or a translation route.
\cite[pp.~254--256]{Street2005LogicSupplement}

Avicenna distinguishes a conjunctive syllogism, whose conclusion
is not explicitly presented in a premise, from a repetitive syllogism,
where the conclusion or its contradictory occurs explicitly within
the premises. This is not simply the division between arguments
using categorical and hypothetical propositions. Averroes criticized
the Avicennan division; he and al-Farabi retained a different division.
Street identifies its relation to perfecting arguments and reduction,
but expressly does not develop the entire proof-theoretic dispute.
\cite[pp.~255, 264, notes~23--24]{Street2005LogicSupplement}

\begin{hbexample}{def:p13-arabic-two-argument-forms}{Two structures, not two interchangeable taxonomies}
In an editorial categorical example, every member of a stipulated
class of triangles is a polygon, and every polygon is a plane figure;
the conclusion that every such triangle is a plane figure is not
explicitly stated in either premise. In a conditional detachment,
the premises say that a specified test's success entails its recorded
signal and that the test succeeds; the signal occurs already as part
of the conditional premise. This illustrates the stated
conjunctive/repetitive criterion. It does not reconstruct all historical
schemata or make every modern conditional one of Avicenna's
historical syllogisms. Terms and conditional retain the same sense
throughout the example.
\end{hbexample}

Formal productivity is distinct from syllogistic matter. In the
Avicennan tradition described by Street, differences in the epistemic
status of premises help distinguish demonstration, dialectic, rhetoric,
poetry and sophistry. A valid pattern does not by itself turn an
accepted opinion into demonstrative knowledge. Later manuals could
emphasize formal questions while retaining a short account of these
differences. Ibn Khaldun's description of this reorganization is a
historical testimony used by Street, not an inventory of every later author.
\cite[pp.~251--252, 255--256]{Street2005LogicSupplement}

Educational reception also involved argument, not automatic adoption.
Al-Farabi analyzed theological and juristic reasoning through syllogisms;
al-Ghazali defended the utility of logic and prefaced his juristic work
\emph{al-Mustasfa} with a logical introduction. Opposition persisted,
including Ibn al-Salah's objection to logic as an avenue to philosophy.
Ibn Taymiyya's criticism of its utility differs from denying the validity
of every syllogistic form. These are positions in the identified historical
discussion, not current religious rulings.
\cite[pp.~252--254, 263, notes~17--21]{Street2005LogicSupplement}

\section{Modal conversion: a real disagreement over an inference}
\label{sec:p13-arabic-modal-conversion}

The problem is not settled by translating one word as \emph{necessary}.
Street begins with an asymmetry of mixed first-figure arguments:
an unmodalized minor combined with a necessary major is accepted
with a necessary conclusion, whereas a necessary minor combined
with an unmodalized major is not. He then asks why a universal
necessary proposition should convert to a particular necessary
proposition. Modern distinctions between necessity attached to the
predication and necessity attached to the whole proposition illuminate
the problem; Aristotle does not state that modern distinction in
the passages being discussed.
\cite[pp.~256--257]{Street2005LogicSupplement}

The Farabian approach discussed here restricts necessary conversion
by distinguishing essential from accidental predication. A literate
being that must be human does not make some human necessarily
literate. An essential predication, such as humanity's relation to
animality in the illustration, is treated differently. Averroes' defense
belongs to a changing exegetical project: Street's note~27 warns
that Averroes changed his position and that this account concerns
his final attempt. It is not an invariant doctrine of every Averroist work.
\cite[pp.~257--258, 264, note~27]{Street2005LogicSupplement}

Avicenna instead rejects unrestricted necessary conversion. His
counterexample concerns a laughing being and a human: the converse
can assert that some human laughs without asserting that a human
must laugh. On this issue he changes a rule, rather than only
excluding a problematic instance from Aristotle's rule. The
Farabian and Avicennan repairs therefore have different mechanisms.
\cite[pp.~258--259]{Street2005LogicSupplement}

\begin{hbexample}{def:p13-arabic-modal-import-case}{A conversion loses its modal strength}
Consider one stipulated human learner who is literate in an initial
state and is not literate in an admissible alternative state, while
remaining human in both. Whenever this individual is literate,
the individual is human; the learner also supplies an actual witness
that a human is literate. But no individual in this one-person model
is literate in every state. Preserving an actual witness for the
converse therefore does not preserve necessity of literacy.
This is a modern separating example for the ambiguity, not an
asserted semantics for Avicenna's full syllogistic. A universal
premise with no actual literate subject would not supply even
that actual particular witness: existence must also be fixed.
\end{hbexample}

\section{Substance and description qualify different time spans}
\label{sec:p13-arabic-substance-description}

Avicenna distinguishes necessity conditioned on the existence of a
substance from necessity conditioned on its being described in a
particular way. The \emph{dhati} reading follows the individual
substance; the \emph{wasfi} reading retains the qualifying description.
Street contrasts a seated person's capacity to stand later with
standing while seated. Similar contrasts arise for a bachelor
becoming married or a sleeper waking. Removing the qualifying
description changes the proposition, not just its wording.
\cite[pp.~259--260]{Street2005LogicSupplement}

\begin{hbexample}{def:p13-arabic-description-time-case}{Possible change is not simultaneous description}
Take a person seated during one interval and standing in a later
interval, and stipulate that sitting and standing exclude each other
at any one instant. What the seated individual can later become
has a witness. Standing during the very interval in which the
individual is seated does not. Retaining the individual while
dropping the time-bound description explains the difference.
The example is not a complete historical modal semantics, nor
a claim about every actual seated person's physical capabilities.
\end{hbexample}

Katibi's later system uses these distinctions without making
Aristotelian exegesis its organizing task. It is based on Avicenna
but modifies his conversion results; Street reports a converse
with a description-conditioned occurrence rather than Avicenna's
one-sided possibility. The shorter chapter omits Katibi's proof
of that conversion. Its report is not a supplied derivation of
the full later modal system.
\cite[pp.~260--262]{Street2005LogicSupplement}

\section{Signification: a lexical route is not a predication}
\label{sec:p13-arabic-signification}

Street's signification study follows a precisely delimited controversy:
Avicenna's \emph{al-Isharat wa-al-Tanbihat}, Pointer~1.6, its
twelfth- and thirteenth-century reception, and Katibi's
\emph{Shamsiyya}. Correspondence concerns the meaning for
which an expression is imposed; containment reaches a constituent
of that meaning; implication reaches a necessarily accompanying
meaning extrinsic to it. Tusi emphasizes that the latter two
combine lexical imposition with an intellectual transition between
meanings. These are signification relations, not interchangeable
ways of affirming a predicate of an object.
\cite[pp.~101--108, Appendix~1]{Street2020PointersSignification}

\begin{hbexample}{def:p13-arabic-signification-case}{Triangle, roof and a human capacity}
In the supplied interpretation, the word for a triangle is imposed
for a three-sided figure. It can thereby convey the constituent
meaning of figure without being independently imposed as the word
for every figure. Avicenna's roof-and-wall illustration instead
uses an extrinsic accompaniment. Conveying wall through roof does
not license the predication that a roof is a wall. The other
illustration concerns a human's receptivity to writing, not that
every human actually writes. Tusi uses these differences to
distinguish predicable from nonpredicable implicates and capacity
from exercise. The illustrations presume their stated account
of accompaniment; they are not new engineering or psychological laws.
\cite[pp.~105--108, 124--125]{Street2020PointersSignification}
\end{hbexample}

\section{Equivocation, evident implication and scientific definition}
\label{sec:p13-arabic-postavicennan-repair}

An expression can receive more than one imposition. The expression
for possibility can name two-sided or one-sided possibility; the
second meaning is also a constituent of the first. Likewise an
expression for the sun can name a celestial body or, by another
imposition, its light. Listing expression and resulting meaning
does not establish the route as correspondence, containment or
implication. Razi's criticism targets that ambiguity. Tusi adds
an intellectual-transition condition; Khunaji defines indirect
significations relative to the imposition through which the first
meaning is conveyed, and Katibi adopts that formulation. Street
reads these as related solutions with distinct argumentative roles.
\cite[pp.~109--112]{Street2020PointersSignification}

\begin{hbexample}{def:p13-arabic-equivocation-case}{Keep the imposition in the inference record}
Under one convention, possible means that an assertion may hold
and may fail. Removing the second component yields a weaker
one-sided condition. Under another convention the same expression
directly names that weaker condition. The weaker content may
be obtained in both cases, but one route uses a constituent of
a previously signified meaning and the other uses a separate
lexical assignment. An argument that silently changes route
cannot establish that the technical uses are identical. This
restates the source's equivocation problem, not a new modal calculus.
\end{hbexample}

Khunaji and Katibi distinguish two strengths of mental implication.
Conceiving both meanings can suffice to recognize their connection
without conceiving the first automatically bringing the second
to mind. The latter is the stronger condition used for implicational
signification. Understanding father involves the correlative child;
finding many truths that follow from a concept does not show
that its expression conveys all of them. An argument about
indefinitely many weak consequences therefore does not immediately
establish indefinitely many signified meanings.
\cite[pp.~113--115]{Street2020PointersSignification}

Containment and implication presuppose correspondence through
which their base meaning is conveyed. Reverse implications
require more: a correspondent meaning need not have proper
conceptual parts, and existence of a consequence does not
establish strong mental implication. Street records disagreement
between Katibi's \emph{Shamsiyya}, his \emph{Jami al-Daqaiq}
and later commentary. A phrase in one work is not promoted
to an invariant doctrine. Grasping a part through its role as
a part differs from additionally conceiving the explicit
qualification of being a part.
\cite[pp.~115--119]{Street2020PointersSignification}

Razi's objection to implication in the sciences concerns precision
in conveying what a thing is. It is not a ban on predicates
extrinsic to a subject's definition: a predicate can itself be
named by correspondence. Nor does Razi's argument simply
repeat Ghazali's infinity objection. Street distinguishes them
and notes where Tusi's report runs them together. Tusi allows
descriptions and incomplete definitions in some contexts while
excluding them as substitutes for complete answers about
quiddity. The scope of Razi's proposed exclusion remains an
interpretive qualification in Street's account.
\cite[pp.~119--123, 125]{Street2020PointersSignification}

The evidence is specific. Street uses Katibi's \emph{Jami
al-Daqaiq} in British Library Or.~11201/2, including folios~11r,
12v and 13r; his first appendix compares Tusi's text with
Leiden Or.~95 and Yeni Cami~763. His second appendix compares
the \emph{Shamsiyya} passage with al-Astana al-Radawiyya~1114,
the \emph{Safina-yi Tabriz} and Trinity R.13.54, rejecting the
last catalogue's holograph attribution. These are Street's
witnesses and conclusions, not manuscripts freshly inspected
here or relationships inferred from shared words alone.
\cite[pp.~102--103, 111--112, 118--121, 124--126, notes~12--13,
37, 61, 64, 73, 78 and 80]{Street2020PointersSignification}

\begin{proofboundary}{def:p13-arabic-source-boundary}{Historical reconstruction and formal scope}
These examples explain particular differences in inferential
organization, modal import and signification. They do not supply
every mood of Avicenna's syllogistic, every successive modal
position of Averroes, or Katibi's complete conversion proof.
The shorter 2005 chapter and bounded signification controversy are
distinct from Street's 2004 historical presentation in
Chapter~\ref{sec:p13-street-arabic-logic-2004}. A
historical calculus requires its selected text, hypothesis and
existence conventions, and all proof and counterexample cases.
A modern separating model is useful at its declared scope,
but is not a verbatim historical semantics or transmission evidence.
\end{proofboundary}
